\documentclass[10pt]{article}
\usepackage{amsmath,amssymb,geometry}
\usepackage[hidelinks]{hyperref}
\geometry{margin=1in}
\title{Snort is not the negative of Col: a counterexample to Conjecture 00000008887}
\author{gaochengzhecpu}
\date{}
\begin{document}
\maketitle

Take the graph $K_2$, with two vertices joined by one edge. In Snort, the
first player wins. In Col, the second player wins. Negating Col exchanges
the two players' options, so the second player still wins. Therefore
\[\operatorname{Snort}(K_2)\ne-\operatorname{Col}(K_2).\]
The graph has a nonidentity automorphism exchanging its two vertices, so
it is not covered by the source's proposed exception for asymmetric graphs.

\paragraph{Rules and complete play analysis.}
Both games use normal play: the player unable to move loses. Each player
places their own color on an empty vertex. In Snort, adjacent vertices may
not have opposite colors; in Col, adjacent vertices may not have the same
color. These are the standard rules; see Definition 1.4 in
\href{https://library.slmath.org/books/Book71/files/GONC6-18.pdf}{Games of No Chance 6}.

In Snort on $K_2$, the first player colors either vertex. The remaining
vertex is adjacent to the opponent's color, so the next player has no
legal move. This is a first-player win, irrespective of who starts.
In Col, after the first move the other player can color the other vertex:
it is adjacent to the opposite color, which is permitted. No empty vertex
remains, so the starting player loses. These cases exhaust every possible
play from the empty board. Equivalently, the game forms are
\[\operatorname{Snort}(K_2)=\{1\mid-1\},\qquad
  \operatorname{Col}(K_2)=\{-1\mid1\}=0.\]
The outcome separation alone suffices; the proof does not require the
numeric simplification of the second form.

\paragraph{Formal verification.}
The Lean file defines a finite normal-play partisan game by its complete
state list, all Left and Right moves, and a strictly decreasing integer
height. Its recursive winning predicate examines every legal move and
declares a terminal player losing. A separate finite-depth evaluator is
proved equal to this well-founded predicate whenever its depth exceeds
the position height; the computation is therefore not an unverified
truncation of play.

The concrete board is a pair of colors in $\{0,1,2\}$, giving all nine
colorings. The file proves that every function from the two vertices to
the three colors is represented. Every move is proved to be exactly a
legal placement at one of the two vertices, and every such move strictly
decreases the number of empty vertices. The adjacency relation, absence
of loops, and nonidentity vertex-exchange automorphism are also checked.

Negation interchanges Left and Right moves at every state. Disjunctive
addition is defined for arbitrary finite games. The predicate
\texttt{EqualValue} expresses equality of outcomes after adding every
finite normal-play context, for either starting player. This is the
standard observational definition of equality of short-game values.
The theorem \texttt{conjecture8887\_counterexample} disproves it for Snort
and negative Col by specializing to the zero-game context, where their
outcomes are respectively winning and losing. Thus the formal conclusion
compares actual games, including the negation operation, rather than
merely comparing names or assigned numerical labels.

The proof uses Lean 4.19.0 with \texttt{Std} only. All finite computations
are checked by the kernel; there are no admitted lemmas, native evaluation
proofs, or custom axioms.
\end{document}
